% 大爆炸（综述）
% license CCBYSA3
% type Wiki

本文根据 CC-BY-SA 协议转载翻译自维基百科\href{https://en.wikipedia.org/wiki/Big_Bang}{相关文章}。

\begin{figure}[ht]
\centering
\includegraphics[width=8cm]{./figures/7a2c26b638cb807a.png}
\caption{宇宙膨胀的时间线：其中空间在每个阶段以圆形剖面示意表示。左侧展示了宇宙暴胀时期的剧烈膨胀；中间部分则表示宇宙膨胀的加速阶段（艺术示意图；时间与尺度均非按比例绘制）。} \label{fig_BBang_1}
\end{figure}
“大爆炸”（The Big Bang）是一种物理理论，用以描述宇宙如何从最初高密度、高温的状态中膨胀而来。基于大爆炸概念的各种宇宙学模型能够解释广泛的天文现象，[1][2][3] 包括轻元素的丰度、宇宙微波背景辐射（CMB），以及宇宙的大尺度结构。宇宙的均匀性——即所谓的视界与平坦性问题——通过“宇宙暴胀”得到解释：那是一段发生在宇宙最早期的加速膨胀阶段。对宇宙膨胀速率的精确测量将最初奇点的时间推算为约 137.87 ± 0.02 亿年前，这被认为是宇宙的年龄。大量实证证据强烈支持这一“大爆炸事件”，该理论如今已被广泛接受。[4]

根据已知物理定律，将这种宇宙膨胀过程反向外推，模型显示宇宙早期处于极端炽热且高密度的原始状态。物理学目前仍缺乏一种被广泛接受的理论来完整描述大爆炸最初阶段的条件。[5] 随着宇宙膨胀，它逐渐冷却到足以形成亚原子粒子，随后又形成原子。这些原初元素——主要是氢，辅以少量氦和锂——在暗物质的引力作用下聚集，最终形成早期的恒星与星系。对超新星红移的观测表明，宇宙的膨胀正在加速，这一现象被归因于一种被称为“暗能量”的概念。

“宇宙膨胀”的思想最早由物理学家亚历山大·弗里德曼（Alexander Friedmann）于 1922 年提出，他在数学上推导出了“弗里德曼方程”。[6][7][8][9] 宇宙膨胀的最早观测依据被称为“哈勃定律”（Hubble’s Law），由物理学家爱德文·哈勃（Edwin Hubble）于 1929 年发表，他发现星系正以与其距离成正比的速率远离地球。独立于弗里德曼的理论工作与哈勃的观测，物理学家乔治·勒梅特（Georges Lemaître）于 1931 年提出，宇宙起源于一个“原初原子”（primeval atom），从而引入了现代意义上的“大爆炸”概念。1964 年，人们发现了宇宙微波背景辐射。随后几年中，对该辐射的观测表明：它在天空各方向上高度均匀，且能量—强度曲线的形状都与“大爆炸模型”预测的早期高温高密度状态相一致。至 20 世纪 60 年代末，大多数宇宙学家已确信竞争性的“稳恒态宇宙”模型是错误的。[10]

然而，观测宇宙中仍存在若干现象尚未能由“大爆炸模型”充分解释。其中包括：物质与反物质数量不对称（即重子不对称性）、星系周围暗物质的具体性质，以及“暗能量”的起源。[11]
\subsection{模型的特征}
\subsubsection{基本假设}
大爆炸宇宙学模型依赖于三项主要假设：\emph{物理定律的普适性}、\emph{宇宙学原理}，以及\emph{物质可被视作理想流体}。[12]  
物理定律的普适性是相对论的基本原则之一。宇宙学原理指出，在大尺度上，宇宙是均匀且各向同性的——无论从哪个方向观察、位于何处，宇宙都呈现相同的总体结构。[13]  
所谓“理想流体”，是指没有黏度的流体，其压力与密度成正比。[14]: 49 

这些假设最初被视为公设，但后来科学家努力对其逐一加以验证。例如，第一个假设（物理定律的普适性）已通过观测得到检验：在宇宙大部分存在时期内，精细结构常数（fine-structure constant）的最大可能偏差量级仅为 $10^{-5}$。[15]  
支撑这些模型的关键物理定律——广义相对论——已在太阳系与双星系统尺度上通过了严格的检验。[16][17]  
宇宙学原理亦得到高度验证：通过对宇宙微波背景（CMB）温度的观测，其一致性水平达到 $10^{-5}$。截至 1995 年，对 CMB 视界尺度的测量表明，宇宙的非均匀性上限约为 10\%。[18]
\subsubsection{膨胀预测}
宇宙学原理极大地简化了广义相对论方程，得出了用于描述宇宙几何结构的 \emph{弗里德曼–勒梅特–罗伯逊–沃尔克度规}（FLRW metric）；结合理想流体假设，可推导出 \emph{弗里德曼方程}，用于描述宇宙几何随时间的演化。  
在这一层次的描述中，唯一的自由参数是质量-能量密度：宇宙的几何形态与其膨胀行为均由密度直接决定。[19]: 73 大爆炸宇宙学的所有主要特征都与这些结果密切相关。[14]: 49 
\subs\begin{figure}[ht]
\centering
\includegraphics[width=6cm]{./figures/818e9d5df373b8f0.png}
\caption{} \label{fig_BBang_2}
\end{figure}

